\documentclass{article}
\usepackage[T1]{fontenc}
\usepackage[utf8]{inputenc}
\usepackage{lmodern}
\usepackage{geometry}
\usepackage{listings}
\usepackage{xcolor}

\geometry{margin=1in}

\definecolor{codebg}{RGB}{245,245,245}
\definecolor{codeborder}{RGB}{220,220,220}

\lstdefinestyle{python}{
  language=Python,
  backgroundcolor=\color{codebg},
  frame=single,
  rulecolor=\color{codeborder},
  basicstyle=\ttfamily\small,
  keywordstyle=\color{blue},
  commentstyle=\color{teal},
  stringstyle=\color{brown},
  showstringspaces=false,
  numbers=left,
  numberstyle=\tiny,
  stepnumber=1,
  numbersep=8pt,
  breaklines=true
}

\title{Sample Python Code in \LaTeX}
\author{}
\date{}

\begin{document}

\maketitle

This short document demonstrates embedding Python code using the \texttt{listings} package.

\section*{Example}

The following snippet defines a simple function and prints a few values:

\begin{lstlisting}[style=python, caption={Simple Python sample}]
import math

def circle_metrics(radius: float) -> tuple[float, float]:
    """Return circumference and area for a circle with the given radius."""
    circumference = 2 * math.pi * radius
    area = math.pi * radius**2
    return circumference, area

for r in (1, 2.5, 5):
    c, a = circle_metrics(r)
    print(f"r={r:.1f} -> circumference={c:.3f}, area={a:.3f}")
\end{lstlisting}

\end{document}
